% !TEX root = ../main.tex
\chapter{Tests, qualité et déploiement}\label{ch:tests}

\section*{Introduction}
Ce chapitre présente la stratégie de test, la chaîne d'intégration et de livraison continues, puis l'architecture de déploiement de la plateforme.

\section{Stratégie de test}
\begin{table}[H]\centering
\caption{Niveaux de test}\label{tab:tests}
\begin{tabularx}{\linewidth}{llX}
\toprule
\textbf{Niveau} & \textbf{Outils} & \textbf{Périmètre} \\
\midrule
Unitaire & JUnit 5, Mockito, Vitest & Règles métier : génération des créneaux, chevauchements, annulation. \\
Intégration & Testcontainers & Dépôts JPA sur un vrai PostgreSQL, isolation des cliniques. \\
Bout en bout & Playwright & Parcours patient complet : recherche, réservation, annulation. \\
Charge & k6 & 200 utilisateurs simultanés sur la recherche de créneaux. \\
\bottomrule
\end{tabularx}
\end{table}

\section{Résultats}
\begin{table}[H]\centering
\caption{Résultats des tests (valeurs d'exemple)}\label{tab:resultats}
\begin{tabularx}{\linewidth}{Xll}
\toprule
\textbf{Indicateur} & \textbf{Objectif} & \textbf{Mesure} \\
\midrule
Couverture du code back-end & $\geq$ 80 \% & 84 \% \\
Tests bout en bout réussis & 100 \% & 100 \% (36/36) \\
Temps de réponse (95\textsuperscript{e} centile) & < 500 ms & 310 ms \\
Vulnérabilités critiques (Trivy) & 0 & 0 \\
\bottomrule
\end{tabularx}
\end{table}
\todo{Remplacer par vos mesures réelles.}

\section{Intégration et livraison continues}
La figure~\ref{fig:cicd} présente la chaîne CI/CD. Chaque \textit{pull request} déclenche la compilation, les tests et l'analyse de qualité ; la fusion sur la branche principale construit les images Docker et les déploie sur le cluster.

\insererfigure{pipeline-ci-cd.png}{Chaîne d'intégration et de livraison continues}{fig:cicd}

\section{Déploiement}
\insererfigure{uml/06-deploiement.png}{Diagramme de déploiement}{fig:deploiement}

\section*{Conclusion}
Ce chapitre a présenté la stratégie de test, les résultats obtenus, la chaîne CI/CD et l'architecture de déploiement de \nomprojet.
